\section{Methodology}
\label{sec:methodology}

\subsection{Collaborative Filtering}

Collaborative filtering is a family of recommendation techniques based on the premise that users who have 
agreed in the past will agree in the future. The core insight is that user preferences are often correlated: 
if user $u_1$ and user $u_2$ have rated several movies similarly, then we can use movies rated by $u_1$ to 
recommend movies to $u_2$.

Formally, the assumption is:

\[
  \text{If } \text{sim}(u_i, u_j) \text{ is high, then unknown ratings by user } u_i 
  \text{ can be inferred from user } u_j \text{'s ratings.}
\]

\subsection{Memory-Based vs.~Model-Based Approaches}

Collaborative filtering is typically categorized into two classes:

\textbf{Memory-Based Approaches:}
\begin{itemize}
  \item Store the full rating matrix in memory.
  \item At prediction time, compute similarity between the target user and all other users.
  \item Predict by averaging ratings from similar users (weighted by similarity).
  \item Advantages: Interpretable, fast for small datasets, requires no training.
  \item Disadvantages: Scales poorly ($O(|\mathcal{U}|^2)$ per prediction), cannot handle sparsity well.
\end{itemize}

\textbf{Model-Based Approaches:}
\begin{itemize}
  \item Learn a compact representation of users and items.
  \item During training, optimize parameters to reconstruct observed ratings.
  \item Predictions are computed using the learned model, not raw data.
  \item Advantages: Scalable ($O(k)$ per prediction, where $k$ is the latent dimension), handles sparsity 
  better, supports extrapolation.
  \item Disadvantages: Requires training, less interpretable, prone to overfitting.
\end{itemize}

This project employs a model-based approach via matrix factorization.

\subsection{Matrix Factorization and SVD}

Matrix factorization is a model-based approach that decomposes the rating matrix $R$ into a product of 
lower-rank matrices:

\[
  R \approx U \Sigma V^T
\]

where:
\begin{itemize}
  \item $U \in \mathbb{R}^{|\mathcal{U}| \times k}$ is the user factor matrix (low-rank approximation of user preferences).
  \item $\Sigma \in \mathbb{R}^{k \times k}$ is a diagonal matrix of singular values (importance of each latent factor).
  \item $V^T \in \mathbb{R}^{k \times |\mathcal{I}|}$ is the transpose of the item factor matrix (low-rank approximation of item characteristics).
  \item $k$ is the rank or number of latent factors (typically $k \ll \min(|\mathcal{U}|, |\mathcal{I}|)$).
\end{itemize}

The prediction for user $u$ on item $i$ is:

\[
  \hat{r}_{u,i} = U_u \Sigma V_i^T
\]

where $U_u$ is the $u$-th row of $U$ and $V_i$ is the $i$-th column of $V$ (or equivalently, the $i$-th row of $V^T$).

\subsection{Rationale for Matrix Factorization}

Matrix factorization is chosen for this project because:

\begin{enumerate}
  \item \textbf{Handles sparsity}: By learning latent factors, the model can predict ratings for 
  unobserved user-item pairs.
  
  \item \textbf{Computational efficiency}: Inference is $O(k)$, enabling real-time predictions.
  
  \item \textbf{Scalability}: Training can be parallelized; the model size grows with $k$, not dataset size.
  
  \item \textbf{Theoretical foundation}: SVD has well-understood mathematical properties and convergence guarantees 
  under standard assumptions.
  
  \item \textbf{Proven empirical performance}: Matrix factorization achieves state-of-the-art results on 
  benchmark datasets when properly tuned.
  
  \item \textbf{Library support}: The \code{Surprise} library provides robust implementations.
\end{enumerate}

\subsection{Optimization Objective}

During training, the model minimizes a loss function that encourages the reconstructed matrix to match 
observed ratings while avoiding overfitting. A standard formulation is:

\[
  \min_{U, V} \sum_{(u,i) \in \mathcal{R}} (r_{u,i} - U_u V_i^T)^2 + \lambda_U \|U\|_F^2 + \lambda_V \|V\|_F^2
\]

where:
\begin{itemize}
  \item The first term is the reconstruction error (mean squared error on observed ratings).
  \item The second and third terms are $L_2$ regularization penalties on user and item factors.
  \item $\|.\|_F$ denotes the Frobenius norm.
  \item $\lambda_U$ and $\lambda_V$ are regularization hyperparameters controlling overfitting.
\end{itemize}

This objective is optimized using stochastic gradient descent (SGD) or alternating least squares (ALS). 
The \code{Surprise} library uses SGD by default, with configurable learning rates and regularization strengths.

\subsection{Why Not Deep Learning?}

While neural networks (autoencoders, NCF) are state-of-the-art for recommendation, this project uses SVD 
because:

\begin{enumerate}
  \item \textbf{Interpretability}: SVD factors have meaningful structure and can be analyzed.
  \item \textbf{Training efficiency}: SVD converges quickly on modest hardware.
  \item \textbf{Simplicity}: SVD is well-understood and implemented reliably.
  \item \textbf{Laboratory setting}: This project prioritizes breadth (full ML product cycle) over 
  cutting-edge model performance.
\end{enumerate}

In production, migrating to neural approaches would be a straightforward extension.
